\section{Граф}

Граф G је уређени пар (V,E) који се састоји од скупа V чије елементе називамо чворовима, и скупа E чији су елементи парови чворов из V, које називамо гранама. Кажемо да грана спаја два чвора којима је одређена, а за чворове те гране кажемо да су суседни. За грану која спаја два чвора u и v кажемо да је инцидентна са чворовима u и v, а чворове u и v називамо крајевима те гране (u,v). За гране инцидентне са истим чвором кажемо да су суседне. На слици \ref{fig:primer_grafa} је дат пример графа чији су чворови елементи скупа V = \{A,B,C,D,E\}, а гране елементи скупа E = \{(A,B), (A,C), (B,C), (C,D), (D,E)\}. \cite{Stanic2020}

\begin{figure}[ht]
\centering
\begingroup
\shorthandoff{"}
\digraph[scale=0.6]{graphexample}{
    edge [dir=none, color=black];
    node [shape=circle];
    A -> B;
    A -> C;
    B -> C;
    C -> D;
    D -> E;
}
\endgroup
\caption{Неусмерен граф са пет чворова и пет грана}
\label{fig:primer_grafa}
\end{figure}

Уколико гранама доделимо усмерење такав граф називамо усмерени граф, а у супротном за граф кажемо да је неусмерен (прост). Додатно гранама можемо доделити тежину, тако добијамо тежински граф. Граф чије гране немају тежину називамо нетежинским графом. Нетежински графови се могу посматрати као тежински графови чије гране имају једнаку тежину (најчешће 1). У претходном примеру, усмерење грана није назначено, али се сматра да су гране усмерене од првог ка другом чвору у пару, па је тако грана (A,B) усмерена од A ка B. Уколико би се сматрало да су гране неусмерене, онда би грана (A,B) била и усмерена од A ка B и од B ка A. \cite{Stanic2020}

\begin{figure}[ht]
\centering
\begingroup
\shorthandoff{"}
\digraph[scale=0.6]{dirgraphexample}{
    edge [color=black];
    node [shape=circle];
    A -> B;
    A -> C;
    B -> C;
    C -> D;
    D -> E;
}
\endgroup
\caption{Усмерен граф са пет чворова и пет грана}
\label{fig:primer_usmerenog_grafa}
\end{figure}

Степен чвора је број грана инцидентих са датим чвором. У неусмереном графу степен чвора је једноставно број грана које су инцидентне са тим чвором, док у усмереном графу разликујемо улазни степен (broj грана које улазе у тај чвор) и излазни степен (broj грана које излазе из тог чвора). \cite{Stanic2020}

Редак граф је граф чији је број суседа чворова ограничен неком (малом) константом. \cite{vandergrinten_et_al:LIPIcs.SEA.2022.11}


У даљем тексту подразумеваћемо да се ради о усмереним, тежинским, ретким графовима, осим ако није другачије назначено.  

\section{Динамички граф}
\subsection{Врсте динамичких графова}
